\section{An original-ground reference and one common trimming radius}
\label{sec:original_ground_reference}

Lemma~\ref{cor:original_ground_reference_trimming_profile} constructs
one reference from the original ground space and chooses one radius for all
interval trims. Its loss bound retains the whole exterior and is independent
of the candidate dimension. These are auxiliaries to the reference argument
in \cite[lines 1801--1805]{Arad2016RigorousRG}; the polynomial dependence
below is at fixed local dimension and gap. The weighted loss is the corrected
collective-trimming estimate in~\cite{gap:alvv16_trimming_constant}.

\begin{lemma}[A reference including flat spectra]
    \label{lem:original_ground_reference}
    \lean{TNLean.MPS.AreaLaw.exists_physical_ground_isOpenMPS_reference_subspace}
    \leanok
    There is a universal $K>0$ with the following property. Let $q\ge1$,
    $\gamma>0$, and let $H=\sum_i h_i$ on $N$ sites, where
    $0\le h_i\le I$ and $h_i$ is supported on sites $i,i+1$.
    Suppose $\mu\in\operatorname{spec}H$ and every spectral value $z$
    satisfies $z=\mu$ or $z\ge\mu+\gamma$. Put
    $G=\ker(H-\mu I)$, $r=\dim G$, $g=\min(\gamma,1)$, and
    $u=1+\log q+\log(1/g)$. For $0<\varepsilon\le1$, define
    $b=16+\log(1/g)+\log(32r(N+1)/\varepsilon)$.
    There are an integer $B\ge1$ and a subspace $U$ of dimension $r$ such
    that $U$ and $G$ are mutually $\varepsilon$-close, every vector of $U$
    has matrix product bond at most $r^2B$, and
    \begin{align}
        \log B\le K\bigl(u^3/g+g^{-1/4}b^{3/4}u\bigr)(1+\log b).
        \label{eq:original_reference_log_profile}
    \end{align}
    No distinct excited spectral value is required. The gap convention
    in this estimate is recorded in
    \cite{gap:alvv16_approximation_parameters}.
\end{lemma}
\begin{proof}\leanok
    \uses{thm:physical_dg_mps_reference_subspace}
    If a spectral value exceeds $\mu$, apply the physical reference
    construction. Otherwise the gap condition forces every eigenvalue
    to equal $\mu$, so $G$ is the whole Hilbert space. Take $U=G$ and
    $B=1$. Every coefficient matrix has rank at most $q^N=r\le r^2$,
    hence every vector has the required exact matrix product
    representation. All terms on the right of
    (\ref{eq:original_reference_log_profile}) are nonnegative.
\end{proof}

\begin{lemma}[A polynomial bound for the reference rank]
    \label{lem:original_ground_reference_polynomial}
    \lean{TNLean.MPS.AreaLaw.exists_physical_ground_reference_polynomial}
    \leanok
    Fix $q\ge1$ and $\gamma>0$. There are $C>0$ and an integer $P\ge0$,
    before every chain and Hamiltonian satisfying
    Lemma~\ref{lem:original_ground_reference}, such that its reference can
    be chosen at error $\varepsilon=10^{-5}$ with
    \begin{align}
        d=r^2B,\qquad b_{\rm ref}=d^2\le C Z^P,
        \qquad Z=(N+1)(r+1).
        \label{eq:original_reference_polynomial}
    \end{align}
    Its dimension, mutual closeness, and all-vector bond bound $d$ are
    retained.
\end{lemma}
\begin{proof}\leanok
    \uses{lem:original_ground_reference}
    At the fixed accuracy, $b\le b_0+\log Z$, where
    $b_0=16+\log(1/g)+\log(3200000)$ depends only on the gap.
    Put $a_0=u^3/g$ and $a_1=g^{-1/4}u$. The logarithmic reference
    estimate gives
    \begin{align*}
        B&\le C_0 Z^{p_0},
        &C_0&=\exp\bigl(K(a_0+5a_1)b_0\bigr),
        &p_0&=\lceil K(a_0+5a_1)\rceil.
    \end{align*}
    Since $r\le Z$, one may take $C=1+C_0^2$ and $P=4+2p_0$.
    These choices also cover the exact flat-spectrum reference.
\end{proof}

\begin{lemma}[One reference and radius for all interval trims]
    \label{cor:original_ground_reference_trimming_profile}
    \lean{TNLean.MPS.AreaLaw.exists_physical_ground_reference_trimming_profile}
    \leanok
    Retain the hypotheses and constants of
    Lemma~\ref{lem:original_ground_reference_polynomial}. The same
    reference $U$ and integer $B$ can be chosen before every interval,
    candidate space, and sample. Define
    \begin{align}
        \xi=\frac{1}{1000d\sqrt{N+1}}>0,
        \qquad d=r^2B,\qquad b_{\rm ref}=d^2.
        \label{eq:original_reference_radius}
    \end{align}
    For every interval $J$, map $U$ and $G$ to
    $\mathcal H_J\otimes\mathcal H_{J^c}$. They remain mutually
    $10^{-5}$-close, and every vector of the mapped reference has Schmidt
    rank at most $b_{\rm ref}$ across every internal interval cut with
    the whole exterior retained. For every integer $k_{\rm cuts}\le N$,
    \begin{align}
        \xi\sqrt{k_{\rm cuts}b_{\rm ref}}&\le10^{-3},\notag\\
        \xi^{-2}&\le10^6 C Z^{P+1}.
        \label{eq:original_reference_radius_budget}
    \end{align}
    For each separately bounded candidate count $M$, the numerical
    ceiling $B_{\rm out}(M)=10^6b_{\rm ref}(N+1)M^2$ satisfies
    \begin{align}
        B_{\rm out}(M)\le10^6 C Z^{P+1}M^2.
        \label{eq:original_reference_output_profile}
    \end{align}
    The loss in (\ref{eq:original_reference_radius_budget}) has no factor
    depending on $M$. Empty intervals and intervals without internal
    cuts are included.
\end{lemma}
\begin{proof}\leanok
    \uses{lem:original_ground_reference_polynomial,
        lem:alvv-interval-reference-weighted}
    An internal interval prefix has at most two boundaries against the
    full chain. Its reference rank is therefore at most $d^2$.
    The interval/exterior coefficient isometry preserves both closeness
    estimates. Since $k_{\rm cuts}\le N$, substitution into
    (\ref{eq:original_reference_radius}) gives the first bound in
    (\ref{eq:original_reference_radius_budget}). Finally
    $\xi^{-2}=10^6b_{\rm ref}(N+1)$ and $N+1\le Z$ give the remaining
    polynomial bounds.
\end{proof}
